% BHU PYQ -- Professional Exam Template
\documentclass[11pt,a4paper]{article}

\usepackage[a4paper,top=2.5cm,bottom=2.5cm,left=2.8cm,right=2.8cm,headheight=14pt]{geometry}
\usepackage{amsmath,amssymb}
\usepackage[T1]{fontenc}
\usepackage[utf8]{inputenc}
\usepackage{lmodern}
\usepackage{microtype}
\usepackage{fancyhdr}
\usepackage{enumitem}
\usepackage{listings}
\usepackage{hyperref}
\usepackage{lastpage}
\usepackage{parskip}

\hypersetup{colorlinks=true,urlcolor=black,linkcolor=black,
  pdftitle={CSB-601: Database Management System -- BHU PYQ}}

\pagestyle{fancy}
\fancyhf{}
\renewcommand{\headrulewidth}{0.4pt}
\renewcommand{\footrulewidth}{0.4pt}
\fancyhead[L]{\small   \textit{Banaras Hindu University}}
\fancyhead[R]{\small   \textit{CSB-601: Database Management System}}
\fancyfoot[C]{\small Page  \thepage\ of \pageref{LastPage}}

\lstset{
  basicstyle=\small tfamily,
  frame=single,
  breaklines=true,
  columns=flexible,
  keepspaces=true
}


\newlist{parts}{enumerate}{2}
\setlist[parts,1]{label=(\alph*),leftmargin=*,itemsep=4pt,topsep=3pt}
\setlist[parts,2]{label=(\roman*),leftmargin=*,itemsep=2pt,topsep=2pt}

\newcommand{\pts}[1]{\hfill   \textit{[#1]}}


\begin{document}

\begin{center}
  {\large\bfseries BANARAS HINDU UNIVERSITY}\\[2pt]
  {
\normalsize B.Sc.\ (Hons.) Semester VI Examination, 2013--14}\\[6pt]
  \rule{0.8\linewidth}{0.5pt}\\[6pt]
  {\large\bfseries Paper: CSB-601 --- Database Management System}\\[4pt]
  {
\normalsize Time: Three Hours \hfill Maximum Marks: 70}
\end{center}

\rule{\linewidth}{0.4pt}

\smallskip



\noindent   \textit{Instructions: Answer any   \textbf{five} questions including Question~1 which is compulsory. Terms and symbols used have their standard meaning.}

\bigskip




\noindent   \textbf{Question 1.} Fill in the blanks:\pts{10}

\begin{parts}
  \item The number of attributes in a relation is called the \underline{\hspace{3cm}} of the relation.
  \item Immunity of the external schemas to changes in the internal schema is referred to as \underline{\hspace{3cm}}.
  \item When one column of a table refers to the values in another column of the same table, it introduces a \underline{\hspace{3cm}} integrity constraint.
  \item The virtual page table never changes during the execution of a transaction --- this property is called \underline{\hspace{3cm}}.
  \item Along with the   \texttt{GROUP BY} clause, one can use the \underline{\hspace{3cm}} clause to define a condition on the group of tuples.
  \item \underline{\hspace{3cm}} and \underline{\hspace{3cm}} are two unary operations in relational algebra.
  \item 4NF is based on the concept of \underline{\hspace{3cm}} dependency.
  \item A schedule $S$ is not conflict serializable if the precedence graph contains \underline{\hspace{3cm}}.
  \item \underline{\hspace{3cm}} and \underline{\hspace{3cm}} are the two desirable properties of relational schema decomposition.
  \item \underline{\hspace{3cm}} is an abstraction through which relationships are treated as higher-level entities.
\end{parts}

\bigskip




\noindent   \textbf{Question 2.}\pts{$3+4+3+5=15$}

\begin{parts}
  \item Explain the difference between primary key and superkey.
  \item What are multi-valued attributes, and how can they be handled within the database design?
  \item What is meant by a recursive relationship type? Give an example.
  \item Define cardinality ratio. Draw an ER diagram to illustrate different cases of cardinality ratio.
\end{parts}

\bigskip




\noindent   \textbf{Question 3.}\pts{$6+6+3=15$}

\begin{parts}
  \item What is the lossless-join property of decomposition? Suppose that we decompose the schema $R = (A, B, C, D, E)$ into $R_1 = (A, B, C)$ and $R_2 = (A, D, E)$. Check whether this decomposition is lossy or lossless if the following set $F$ of functional dependencies holds: $F = \{A  o BC,\ CD  o E,\ B  o D,\ E  o A\}$.
  \item Define full, partial, and transitive functional dependencies. Explain the normal form based on transitive functional dependency with a suitable example.
  \item Define multi-valued functional dependency.
\end{parts}

\bigskip




\noindent   \textbf{Question 4.}\pts{$6+6+3=15$}

\begin{parts}
  \item Discuss the immediate update technique of recovery. How is it different from the deferred update technique? Why is it called the UNDO/REDO method?
  \item What is the two-phase locking protocol? What are some variations of the two-phase locking protocol? Why is strict two-phase locking often preferred?
  \item What is a recoverable schedule? Why is recoverability of schedules desirable?
\end{parts}

\bigskip




\noindent   \textbf{Question 5.}\pts{$5+10=15$}

\begin{parts}
  \item With a suitable example, explain the difference between natural join, equi-join, and Cartesian product operations.
  \item Consider the following relational database. Answer each of the queries in SQL:

  \smallskip
   \texttt{Branch}(\underline{bname}, bcity, assets)\\
   \texttt{Customer}(\underline{cname}, street, ccity)\\
   \texttt{Depositor}(\underline{cname}, \underline{account\#})\\
   \texttt{Account}(\underline{bname}, \underline{account\#}, balance)

  
\begin{enumerate}[label=(\roman*)]
    \item Find account numbers of accounts with balances between Rs.\,70,000 and Rs.\,1,00,000.
    \item Find all customers whose street includes the substring ``main''.
    \item Find branches and their average balances where the average balance is more than Rs.\,2,000.
    \item Delete the records of all accounts with balances below the average.
    \item Increase all balances by 10\% whose balance is greater than the average.
  \end{enumerate}
\end{parts}

\bigskip




\noindent   \textbf{Question 6.}\pts{$6+3+6+5+4+6=30$}

\begin{parts}
  \item Describe the ANSI/SPARC three-level architecture. How are these three levels related to the concept of physical and logical data independence?
  \item Describe the steps of query processing with a flow diagram.
  \item Explain any three heuristic rules for query optimization through examples.
  \item Define weak entity type. Explain how it is handled in relational database designing.
  \item Explain the Select and Project operators in Relational Algebra through an example.
  \item What are the database security issues? What is the difference between discretionary and mandatory access control?
\end{parts}

\vfill
\rule{\linewidth}{0.4pt}

\begin{center}\small
\href{https://bhu-exam.vercel.app/}{Click here to visit our website}\\[4pt]\href{https://me-aryan.vercel.app/}{Click here to visit our contributor}\\[4pt]\href{https://quashan-me.vercel.app/}{Click here to visit our contributor}
\end{center}
\end{document}